\section{Transmisión, ingesta e interfaces}
\label{sec:transmision-interfaces}

Una vez acondicionada y digitalizada, la señal debe abandonar cada nodo, llegar sin
ambigüedades a un punto común, conservar su contexto de adquisición y quedar disponible
para revisión y procesamiento:

\begin{equation}
\text{nodos de medida}
\longrightarrow
\text{maestro de campo}
\longrightarrow
\text{servidor de ingesta}
\longrightarrow
\text{revisión y MASW}.
\end{equation}

La primera interfaz controla la adquisición (armar nodos, ajustar parámetros, iniciar o
detener una captura, observar el enlace, exportar); la segunda controla el conjunto de
datos (validar, filtrar, agrupar, alinear, promediar, invertir). Esta separación evita
que el microcontrolador ejecute tareas pesadas y permite seguir trabajando en campo sin
conexión a Internet.

\subsection{Plataforma de comunicación}

Se eligió la familia ESP32 por accesibilidad local, bajo costo, documentación abundante
e integración de Wi-Fi; no se lo usa como instrumento de medida, sino después de la
digitalización: el nodo esclavo empaqueta los datos del frente analógico y los entrega
al maestro. El enlace entre nodos usa ESP-NOW, que da una red directa de baja sobrecarga
sin depender de un \emph{router} en el sitio, a cambio de carga útil limitada por
paquete y entrega no confiable por sí sola; el diseño compensa esto con fragmentación,
numeración de mensajes, confirmaciones, reintentos y almacenamiento local previo a la
subida. El maestro concentra los enlaces ESP-NOW y expone además una red Wi-Fi propia
con una interfaz web embebida accesible desde cualquier teléfono o computadora sin
instalar software; cuando hay conectividad disponible pasa a una fase de enlace,
se asocia a una red externa y transfiere la cola de capturas pendientes.

\subsection{Por qué se desarrollaron varias interfaces}

Cada interfaz resolvió una pregunta de una etapa distinta y dejó funciones que se
trasladaron a la siguiente versión. Se empezó por un osciloscopio MATLAB de un solo
canal para fijar las variables que debían quedar controladas (puerto, referencia DC,
ganancia, filtro) antes de comparar capturas; al incorporar varios nodos, una segunda
aplicación MATLAB agregó una vista de estado global y una pestaña por esclavo, tratando
ya la adquisición como un sistema y no como un canal aislado. Una versión de escritorio
en PyQt reorganizó esto en una herramienta multi-nodo con temas claro/oscuro, útil tanto
en laboratorio como a la intemperie. El paso siguiente fue eliminar la dependencia de
una computadora con MATLAB o Python instalado: la interfaz se embebió como una \emph{SPA}
servida directamente por el maestro, operable desde cualquier navegador y funcional sin
conexión a Internet, con menús de Captura, Esclavos, Preservados, Enlace y Log. En
paralelo se prototipó en PyQt, y luego en un servidor web, el flujo de posprocesado
completo (filtros, agrupamiento, enfase, promedios, \emph{waterfall}, dispersión e
inversión MASW), migrando finalmente a un servidor FastAPI para tener una dirección
conocida, estado persistente entre campañas y acceso remoto sin reinstalar nada en cada
equipo de campo.

\begin{table}[ht]
\centering
\small
\begin{tabularx}{\textwidth}{@{}p{0.20\textwidth}p{0.23\textwidth}Y@{}}
\toprule
\textbf{Interfaz} & \textbf{Uso principal} & \textbf{Aporte al diseño} \\
\midrule
MATLAB: alcance USB & Verificar el canal individual & Visualización cruda, referencia DC, FIR, guardado de muestras. \\
MATLAB: InterfaceESP & Probar maestro y esclavos & Sincronización, lotes, control de PGA/FIR por nodo, depuración. \\
PyQt: adquisición & Operar varios nodos & Escritorio multipestaña, estado de enlace, gráficos simultáneos. \\
SPA embebida & Operar en campo sin instalar software & Control desde el navegador, funcionamiento fuera de línea. \\
PyQt: revisión & Explorar el flujo científico & Filtros, grupos, enfase, promedios, \emph{waterfall} y MASW. \\
Servidor FastAPI & Ingestar y procesar en una dirección conocida & Acceso remoto, estado persistente entre campañas. \\
\bottomrule
\end{tabularx}
\caption{Funciones y alcance de las interfaces desarrolladas.}
\label{tab:evolucion-interfaces}
\end{table}

\begin{figure}[ht]
\centering
\includegraphics[width=0.8\textwidth,height=0.62\textheight,keepaspectratio]{esp_historico/07_pga_pgaout/01_captura_completa.png}
\caption{Estado vigente de la SPA embebida: cinco menús (Captura, Esclavos,
Preservados, Enlace, Log), decimación, límites de RAM/SD, selección PGA/PGAout y
subida al servidor. Frente a la primera versión ejecutable, la interfaz pasó de
exponer sólo transporte a exponer una operación de campo completa.}
\label{fig:spa-hito7-resumen}
\end{figure}

\subsection{Resultados reproducibles: estado actual y límites}

Además de las interfaces, el repositorio conserva salidas numéricas generadas por los
scripts de análisis sobre 21 trazas reales entre 10 y \SI{50}{\meter}, muestreadas a
\SI{1020}{\hertz}. Comparando el seguimiento automático de cresta previo contra un
seguidor Kalman+RTS, el error respecto de una referencia externa en la banda
8--\SI{29.83}{\hertz} bajó de aproximadamente 138--148 a \SI{16}{\meter\per\second}: es
evidencia de mejora del \emph{picking} en este conjunto, no una validación general del
sistema. Una inversión preliminar con esa curva (25 puntos entre 8 y \SI{21.71}{\hertz})
alcanza un RMSE de \SI{0.463}{\meter\per\second} contra el \emph{pick}, resolviendo la
banda disponible hasta aproximadamente \SI{5.4}{\meter}: la cadena completa
trazas~$\to$~dispersión~$\to$~inversión~$\to$~$V_S(z)$ ya funciona de punta a punta,
pero todavía falta evaluar repetibilidad entre grupos y campañas, sensibilidad a la
geometría, incertidumbre del \emph{picking}, no unicidad de la inversión y comparación
con información independiente del sitio.

\subsection{Acceso remoto y conclusión}

Como el equipo de campo opera sin Internet pero la ingesta remota necesita una dirección
alcanzable, se combina una VPN de malla con tunelamiento para publicar el servidor bajo
una dirección conocida, sin exigir configuración manual de puertos en cada red visitada;
esta capa no reemplaza el respaldo local, que se libera recién después de confirmar la
transferencia. La evolución MATLAB~$\rightarrow$~PyQt~$\rightarrow$~SPA embebida y
servidor web separó responsabilidades ---adquisición local tolerante a desconexiones,
transporte verificable, ingesta persistente y procesamiento reproducible--- en lugar de
concentrar toda la operación en un script único. Quedan pendientes pruebas integrales
con todos los nodos conectados, medición de pérdida y latencia bajo interferencia, y
verificación de recuperación ante cortes de energía o de red.
